% =====================================================================
% draft_evaluation.tex
% Experimental Setup + evaluation scaffolding for
%   AAAI_27_Rodent/AnonymousSubmission2027.tex
%
% Covers revision-plan items A5, A7, A8.2 (Experimental Setup budget),
% and provides the empty result containers for Track B items
% B1-B12.
%
% ---------------------------------------------------------------------
% CONVENTION USED THROUGHOUT THIS FILE -- READ BEFORE EDITING
% ---------------------------------------------------------------------
% There is not one measured number in any table below. Every cell that
% must receive a measurement renders as
%       [R-E]
% via the \REQEXP macro defined immediately below, and every table is
% preceded by a block comment beginning
%       % REQUIRES_EXPERIMENT:
% that names exactly what must be run to fill it.
%
% Do not "temporarily" put a plausible number in a [R-E] cell to see how
% the table looks. A placeholder that reads like a result is the single
% most damaging thing that can happen to this manuscript. If you need to
% check the layout, put 00.0 in EVERY cell of the table at once so that
% no partially-filled version can ever escape.
%
% The four numbers already in the frozen abstract -- 91.5 / 86.8 / 81.2 /
% 83.5 -- appear in exactly one place in this file (the protocol-ladder
% table, P1 row) and are labelled there as the figures produced by the
% originally submitted configuration. They are not repeated elsewhere.
%
% ---------------------------------------------------------------------
% PAGE BUDGET
% ---------------------------------------------------------------------
% Compiled standalone against the AAAI-27 template this file occupies
% roughly 4.5 columns: ~1.5 for Experimental Setup and ~3 for the eight
% tables and their surrounding prose. The revision plan (A8.2) budgets
% 0.70 pg for Experimental Setup and 1.60 pg for Results.
% It therefore OVERRUNS as written, deliberately -- it is a superset, and
% the cut order is given in the block comment that opens the results
% scaffolding below. Cut tables, not the honesty paragraph in the
% protocol subsection.
%
% ---------------------------------------------------------------------
% PACKAGE / PREAMBLE REQUIREMENTS
% ---------------------------------------------------------------------
% This file needs only booktabs and amsmath, both already loaded by the
% manuscript preamble. It requires NO new packages.
% The \REQEXP definition must be moved to the preamble (or kept here, but
% it must appear exactly once in the merged document -- if another draft
% file defines it too, delete one copy).
% ---------------------------------------------------------------------

\newcommand{\REQEXP}{\texttt{[R-E]}}

% =====================================================================
% NOTE ON PLACEMENT
% The revision plan (A8.2) gives Experimental Setup its own section with
% a 0.70-page budget, split out of the current Methodology. That is what
% is written below. If the authors prefer to keep a single Methodology
% section, demote \section -> \subsection and each \subsection ->
% \subsubsection throughout; nothing else changes.
% =====================================================================

\section{Experimental Setup}
\label{sec:setup}

\subsection{Cohort and Provenance}
\label{subsec:cohort}

% REQUIRES_EXPERIMENT: none of this is an experiment -- every blank below
% is a fact the authors already hold and must transcribe. It is flagged
% the same way because writing a guess would be fabrication either way.
% Needed: species (the manuscript never says mouse or rat, while
% comparing itself against mouse- and rat-specific published figures),
% strain, sex, age at recording, housing and light cycle, electrode
% montage and reference scheme, acquisition hardware, animal-use protocol
% and approving body, and the licence under which the recordings may be
% redistributed.
Recordings were obtained from eight
% REQUIRES_EXPERIMENT: species and strain
laboratory rodents monitored continuously for two days, digitized at
$f_s = 500$\,Hz from a single EEG channel and segmented into
non-overlapping $10$\,s epochs, giving $138{,}240$ epochs in total.
% REQUIRES_EXPERIMENT: per-animal epoch counts. Note that the cohort
% total is exactly 8 x 48 h at 10 s resolution, which is consistent with
% no epoch having been excluded from any recording. If artifact rejection
% was in fact applied, the counts will not be uniform and the total will
% not be 138,240 -- state which is the case.
Per-animal epoch counts are given in Table~\ref{tab:classdist}.
No electromyogram, accelerometer, or video channel was available, so
every result reported here is obtained from cortical EEG alone; this
distinguishes the present setting from scorers that consume paired
EEG and EMG \citep{jha2024_slumbernet, wang2023_intellisleepscorer}
and it bounds what any of the reported figures can be compared against.

\subsection{Reference Standard}
\label{subsec:labels}

% REQUIRES_EXPERIMENT: label provenance. The manuscript cites a human
% inter-rater kappa of 0.76-0.85 twice but never states how its own
% ground truth was produced, which means the reported accuracy is
% currently measured against an unspecified target. Needed: number of
% scorers, their training, the scoring criteria and epoch length used,
% whether any recording was double-scored, the observed agreement on the
% double-scored subset, and how disagreements were resolved.
Vigilance-state labels were assigned by
% REQUIRES_EXPERIMENT: number and qualification of scorers
expert scorers following
% REQUIRES_EXPERIMENT: scoring criteria and any published standard followed
established criteria. Because manual scoring of rodent EEG is itself
subject to inter-rater disagreement in the range $\kappa = 0.76$--$0.85$,
agreement with this reference standard is bounded above by the
reliability of the standard itself. We therefore report Cohen's $\kappa$
against the label set alongside accuracy: $\kappa$ is the only quantity
we compute that is on the same scale as the published inter-rater
figures, and it is the appropriate one to compare against them.

\subsection{Evaluation Protocols}
\label{subsec:protocols}

Because the epochs in this dataset are neither independent of one another
nor exchangeable across animals, the choice of partition determines what
a reported figure means. We therefore define three protocols and report
every headline metric under each. They differ only in the unit that is
held out; the feature pipeline, hyperparameters, and random seed are
identical throughout.

\begin{itemize}\itemsep2pt\parskip0pt
    \item \textbf{P1 (epoch-wise).} All $138{,}240$ epochs are pooled
    across animals and partitioned at random, stratified by vigilance
    state, into $80\%$ training and $20\%$ test. This is the protocol
    under which the figures originally reported for this system were
    produced.
    \item \textbf{P2 (contiguous-block, within-animal).} Each recording
    is divided into contiguous blocks of
    % REQUIRES_EXPERIMENT: fix and state the block unit -- whole hours,
    % or whole light/dark phases. Whole light/dark phases are preferable
    % because they also break the circadian confound; whole hours give
    % more folds. Pick one and use it everywhere.
    fixed duration, and whole blocks are assigned to training or test.
    Every animal still contributes to both partitions, but no test epoch
    is temporally adjacent to a training epoch.
    \item \textbf{P3 (leave-one-animal-out).} Eight folds, each holding
    out all epochs of one animal. No epoch from a test animal, and no
    epoch recorded in the same session as a test epoch, is seen during
    training or model selection.
\end{itemize}

\paragraph{What P1 establishes, and what it does not.}
Under P1 every animal contributes epochs to both partitions, and because
the recordings are continuous and the epochs are consecutive, the
immediate temporal neighbours of almost every test epoch---drawn from the
same sleep bout, the same electrode, the same session, and the same
animal---are in the training set. P1 therefore measures how well the
classifier interpolates within recordings it has already partly seen.
That quantity is not meaningless: it is the operating point of the common
laboratory workflow in which a scorer is fitted on manually scored
segments of a cohort's recordings and applied to the remainder of those
same recordings. But it is not a measure of transfer to an unseen animal,
an unseen session, or an unseen laboratory, and we do not present it as
one. Figures obtained under P1 are not comparable to figures published
for scorers evaluated on held-out animals or held-out datasets, and every
comparison in this paper that crosses that line is marked as such.
P3 is the protocol that supports subject-independent claims; P2 separates
the contribution of temporal adjacency from the contribution of subject
identity. We report all three because the difference between them is
itself the measurement, and we report it whichever direction it goes.

% NOTE TO THE AUTHORS: the paragraph above is deliberately flat. Do not
% add a mitigating clause ("nevertheless", "it should be noted that the
% split is standard in this literature", "we argue that"). The value of
% this paragraph to a reviewer is entirely in its refusal to argue. The
% fact that the epoch-wise split is common in this application literature
% belongs in Related Work, stated about other people's papers, not here
% as a defence of this one.

\subsection{Metrics}
\label{subsec:metrics}

Let $\mathcal{Y} = \{\textsc{wake}, \textsc{sws}, \textsc{rem}\}$, let
$y_i$ and $\hat{y}_i$ denote the reference and predicted state of test
epoch $i$, and let $n$ be the number of test epochs. For a state
$c \in \mathcal{Y}$ write $\mathrm{TP}_c$, $\mathrm{FP}_c$, and
$\mathrm{FN}_c$ for the one-vs-rest counts. We report
\begin{align}
\mathrm{Pre}_c &= \frac{\mathrm{TP}_c}{\mathrm{TP}_c + \mathrm{FP}_c},
\qquad
\mathrm{Rec}_c = \frac{\mathrm{TP}_c}{\mathrm{TP}_c + \mathrm{FN}_c},
\label{eq:prerec}\\
\mathrm{F1}_c &= \frac{2\,\mathrm{Pre}_c\,\mathrm{Rec}_c}
                      {\mathrm{Pre}_c + \mathrm{Rec}_c},
\label{eq:f1}
\end{align}
and their unweighted means over the three states, which we label
\emph{macro} throughout. Accuracy is
$\frac{1}{n}\sum_i \mathbf{1}[\hat{y}_i = y_i]$. Balanced accuracy is the
unweighted mean of the per-state recalls and is therefore numerically
identical to macro recall in this three-class setting; we report it once,
under the macro-recall heading, rather than twice under two names.
Chance-corrected agreement is Cohen's
\begin{equation}
\kappa = \frac{p_o - p_e}{1 - p_e},
\label{eq:kappa}
\end{equation}
with $p_o$ the observed agreement and $p_e$ the agreement expected from
the marginal distributions.

% NOTE TO THE AUTHORS: this next sentence exists to close the numeric
% inconsistency identified in the manuscript (Track A item A7): the
% network row currently reports accuracy = recall = 83.54%, which is the
% signature of a WEIGHTED recall reported against a MACRO recall for
% XGBoost. Every metric in every table must be recomputed under one
% stated convention before any of these tables is filled.
Unless a table states otherwise, every precision, recall, and F1 figure
in this paper is macro-averaged, computed with the same convention for
every model and every protocol.

For probabilistic evaluation we report the multiclass Brier score
\begin{equation}
\mathrm{BS} = \frac{1}{n}\sum_{i=1}^{n}\sum_{c \in \mathcal{Y}}
              \bigl(\hat{p}_{ic} - \mathbf{1}[y_i = c]\bigr)^2,
\label{eq:brier}
\end{equation}
which under this convention lies in $[0,2]$, and the expected calibration
error over $M$ equal-width bins $B_1,\dots,B_M$ of the predicted
probability,
\begin{equation}
\mathrm{ECE} = \sum_{m=1}^{M} \frac{|B_m|}{n}
               \bigl|\operatorname{acc}(B_m) - \operatorname{conf}(B_m)\bigr|,
\label{eq:ece}
\end{equation}
together with the maximum calibration error
$\mathrm{MCE} = \max_m |\operatorname{acc}(B_m) - \operatorname{conf}(B_m)|$.
% REQUIRES_EXPERIMENT: state M, and state whether the reliability
% analysis is one-vs-rest per class or top-label. The reliability diagram
% currently in the manuscript shows approximately 20 bins at 0.05
% spacing, while the text says ten -- resolve this from the plotting
% code, not from the image.
Reliability curves are computed
% REQUIRES_EXPERIMENT: one-vs-rest or top-label; number of bins
over equal-width bins of the predicted probability.

\subsection{Baselines}
\label{subsec:baselines}

We compare against four references, chosen so that at least one is
external to this work.
\begin{itemize}\itemsep2pt\parskip0pt
    \item \textbf{Majority-class predictor.} Predicts the most frequent
    state for every epoch. This fixes the floor: on a three-class problem
    with a strongly skewed prior, accuracy alone is uninformative without
    it.
    \item \textbf{Multinomial logistic regression.} $L_2$-regularized,
    fitted on the same design matrix, the same partition, and the same
    standardization as every other model.
    % REQUIRES_EXPERIMENT: regularization strength, solver, maximum
    % iterations, and whether the solver converged. The manuscript's
    % Conclusion reports 54.9% for this model and reports it nowhere
    % else; 54.9% on a three-class problem whose majority class alone
    % plausibly exceeds 50% is close to the trivial rate, and a published
    % logistic-regression baseline on a comparable task reaches 93.3%
    % \citep{wang2023_intellisleepscorer}. Either the figure is
    % reproducible and explicable, or the claim must be deleted.
    \item \textbf{Feed-forward network.} The two-layer network described
    in the Method section, retrained on the \emph{identical} feature
    vector supplied to the boosted ensemble.
    % REQUIRES_EXPERIMENT: the manuscript's network consumed 5,009
    % columns. If the boosted ensemble consumed a different column set,
    % the existing comparison measures inputs rather than models. Retrain
    % both on one matrix or reframe the comparison explicitly as
    % engineered-features-versus-raw-samples.
    \item \textbf{SAGER \citep{saevskiy2025_sensors}.} An open-source
    single-channel rodent scorer combining artifact processing,
    multi-band spectral analysis, and Gaussian-mixture clustering. It is
    the only published rodent scorer we are aware of that is both openly
    available and applicable to EEG without EMG; the widely used
    AccuSleep \citep{barger2019_accusleep} and IntelliSleepScorer
    \citep{wang2023_intellisleepscorer} both require an EMG channel that
    this cohort does not have. We run SAGER on these eight recordings
    under the same protocols and score it against the same labels.
    % REQUIRES_EXPERIMENT: run SAGER (github.com/sykesva/SAGER) on this
    % cohort. Report NO comparison number that was not produced by
    % actually running the method on this data. If integration fails,
    % say so explicitly and report no number at all -- an honest
    % "we could not run it, for these reasons" is publishable; a number
    % transcribed from SAGER's own paper and placed in a table of results
    % on this cohort is not.
\end{itemize}

\subsection{Model Selection and Hyperparameters}
\label{subsec:hparams}

% REQUIRES_EXPERIMENT: the grid actually searched. Needed: each parameter
% name, the exact list of values tried, the resulting number of
% configurations, the number of cross-validation folds, THE UNIT ON WHICH
% FOLDS WERE ASSIGNED (epochs or animals), and the selection criterion.
% Reproducibility-checklist item 4.1 asks for precisely this. If the
% search folds were assigned over pooled epochs, that must be stated,
% because it means the selected hyperparameters saw every animal.
Hyperparameters were selected by grid search over
% REQUIRES_EXPERIMENT: parameter names, value lists, configuration count
a fixed grid, with folds assigned over
% REQUIRES_EXPERIMENT: epochs or animals
and configurations ranked by
% REQUIRES_EXPERIMENT: selection criterion
on the validation folds. The selected configuration uses a learning rate
of $0.1$, $500$ estimators, subsample $0.8$, \texttt{colsample\_bytree}
$0.8$, $\gamma = 0$, $\lambda = 1$, and seed $42$.
% REQUIRES_EXPERIMENT: max_depth is named in the tuning description but
% its selected value is never given; also missing are objective,
% tree_method, min_child_weight, and eval_metric.
The same configuration is used unchanged under P1, P2, and P3, so that
the differences between protocols reflect the partition and nothing else.

% NOTE TO THE AUTHORS -- this is a correctness bug, not a presentation
% issue. The manuscript specifies objective = multi:softmax, which
% returns hard class labels. A reliability diagram, an ECE, and a Brier
% score cannot be computed from hard labels. Either the reported
% objective is wrong and the fitted model actually used multi:softprob,
% or the calibration analysis was computed from a different model than
% the one described. State which, and make the reported objective match
% the model that produced the calibration figure.
Class probabilities required for the calibration analysis are obtained
from
% REQUIRES_EXPERIMENT: state the objective actually used
the multiclass probabilistic objective. No class reweighting or
resampling was applied at any stage; stratification preserves the
population class proportions in each partition by construction and does
not correct for them.
% NOTE: the manuscript currently claims stratification prevented "class
% imbalance issues". It does the opposite -- it guarantees the imbalance
% is reproduced in the test set. The sentence above replaces that claim.
Feature standardization for the models that require it was fitted on the
training partition only and applied unchanged to the test partition.
% REQUIRES_EXPERIMENT: confirm from the code that the scaler was not
% fitted on the full dataset. If it was, that is an additional leakage
% path affecting the network row and must be stated.

\subsection{Compute}
\label{subsec:compute}

% REQUIRES_EXPERIMENT: hardware (CPU model, GPU model, RAM), operating
% system, and versions of xgboost, scikit-learn, numpy, scipy, and
% TensorFlow. Also: wall-clock time for one fit, for the full P3 sweep,
% and for feature extraction over the whole cohort.
A single fit of the selected configuration completes in under four
minutes on
% REQUIRES_EXPERIMENT: hardware
a single GPU. Per-epoch inference latency and throughput are reported in
Table~\ref{tab:compute}, measured on
% REQUIRES_EXPERIMENT: CPU and edge-device specifications
both a workstation CPU and an embedded-class device, since untethered
closed-loop use is the setting in which a compact ensemble has an
advantage over a deep network.


% #####################################################################
% #####################################################################
%  RESULTS SCAFFOLDING
%
%  Everything from here down belongs in the Results section. The tables
%  are ordered by the priority given in the revision plan: cheapest and
%  highest-value first. If the page budget binds, cut from the BOTTOM of
%  this list, in this order:
%     1st to go: Table~\ref{tab:context}   (external context -- optional)
%     2nd to go: Table~\ref{tab:compute}   (edge benchmark -- optional)
%     3rd to go: Table~\ref{tab:calib}     (fold into prose)
%     4th to go: Table~\ref{tab:confusion} (fold into tab:perclass)
%  Do NOT cut tab:perclass, tab:ladder, tab:loao, tab:ablation, or
%  tab:baselines. Each of those answers a question a reviewer will ask
%  in writing.
% #####################################################################
% #####################################################################


% =====================================================================
% TABLE: class distribution
% =====================================================================
% REQUIRES_EXPERIMENT (protocol E2 in EXPERIMENTS.md):
%   Count epochs by (animal, vigilance state) over the whole cohort.
%   No model, no training -- a groupby on the label file. ~10 minutes.
%   Fill: 24 count cells, 3 column totals, 3 percentages, 8 row totals.
% =====================================================================
\begin{table}[t]
\centering
\footnotesize
\setlength{\tabcolsep}{4pt}
\begin{tabular}{@{}lrrrr@{}}
\toprule
\textbf{Animal} & \textbf{Wake} & \textbf{SWS} & \textbf{REM} & \textbf{Total} \\
\midrule
A1 & \REQEXP & \REQEXP & \REQEXP & \REQEXP \\
A2 & \REQEXP & \REQEXP & \REQEXP & \REQEXP \\
A3 & \REQEXP & \REQEXP & \REQEXP & \REQEXP \\
A4 & \REQEXP & \REQEXP & \REQEXP & \REQEXP \\
A5 & \REQEXP & \REQEXP & \REQEXP & \REQEXP \\
A6 & \REQEXP & \REQEXP & \REQEXP & \REQEXP \\
A7 & \REQEXP & \REQEXP & \REQEXP & \REQEXP \\
A8 & \REQEXP & \REQEXP & \REQEXP & \REQEXP \\
\midrule
Cohort & \REQEXP & \REQEXP & \REQEXP & 138{,}240 \\
Share  & \REQEXP & \REQEXP & \REQEXP & 100\% \\
\bottomrule
\end{tabular}
\caption{Epoch counts by animal and vigilance state. Cells marked
\REQEXP\ are unmeasured slots, not results. The cohort row fixes the
class prior against which every accuracy figure in this paper must be
read, and the per-animal rows fix how much that prior varies across the
folds of the leave-one-animal-out protocol.}
\label{tab:classdist}
\end{table}

% ---------------------------------------------------------------------
% PROSE FOR tab:classdist -- states what the table is for, asserts no result
% ---------------------------------------------------------------------
Table~\ref{tab:classdist} gives the composition of the cohort. It is
reported before any performance figure because the three states are not
equally represented: REM occupies a small minority of epochs, so overall
accuracy is dominated by the two majority states and the majority-class
floor in Table~\ref{tab:baselines} is the reference point against which
any accuracy figure must be read. The per-animal rows additionally
determine how far the class prior varies between the folds of P3, which
bears directly on the spread reported in Table~\ref{tab:loao}.


% =====================================================================
% TABLE: per-class metrics
% =====================================================================
% REQUIRES_EXPERIMENT (protocol E1 in EXPERIMENTS.md):
%   Recompute from saved test-set predictions. NO retraining is needed if
%   predictions were saved; if they were not, one refit of the selected
%   configuration (~4 min) regenerates them.
%   Fill: 3 states x 4 columns, the two average rows, accuracy, kappa.
%   Produce one instance of this table per protocol (P1, P2, P3), or add
%   a protocol column and merge.
%
% CONSISTENCY CHECK, to be done the moment this table is filled:
%   Under P1 and the originally submitted configuration, the macro row
%   must reproduce 86.8 / 81.2 / 83.5 and accuracy must be 91.5%. Those
%   are the figures in the frozen abstract. If they do not reproduce,
%   the discrepancy must be found and explained BEFORE anything else is
%   written -- do not adjust the table to match the abstract.
%   If a corrected configuration (e.g. with the subject-identity column
%   removed) produces different values, report BOTH, label each by
%   configuration, and state the relationship in one sentence. Do not
%   silently replace one with the other.
% =====================================================================
\begin{table}[t]
\centering
\small
\begin{tabular}{@{}lcccc@{}}
\toprule
\textbf{State} & \textbf{Precision} & \textbf{Recall} & \textbf{F1} & \textbf{Support} \\
\midrule
Wake & \REQEXP & \REQEXP & \REQEXP & \REQEXP \\
SWS  & \REQEXP & \REQEXP & \REQEXP & \REQEXP \\
REM  & \REQEXP & \REQEXP & \REQEXP & \REQEXP \\
\midrule
Macro avg.    & \REQEXP & \REQEXP & \REQEXP & \REQEXP \\
Weighted avg. & \REQEXP & \REQEXP & \REQEXP & \REQEXP \\
\midrule
Accuracy         & \multicolumn{4}{l}{\REQEXP} \\
Cohen's $\kappa$ & \multicolumn{4}{l}{\REQEXP} \\
\bottomrule
\end{tabular}
\caption{Per-state performance under
% REQUIRES_EXPERIMENT: name the protocol this instance reports (P1, P2,
% or P3) and the configuration (with or without the subject-identity
% column). One table per protocol, or add a protocol column.
protocol~[P\,?]. Cells marked \REQEXP\ are unmeasured slots, not results.
Balanced accuracy equals the macro recall entry and is not listed
separately. All averages are computed under the single convention stated
in the Metrics description.}
\label{tab:perclass}
\end{table}

% ---------------------------------------------------------------------
% PROSE FOR tab:perclass
% ---------------------------------------------------------------------
Table~\ref{tab:perclass} decomposes the aggregate figures by state. The
decomposition is reported because every substantive claim about this task
concerns the minority state: macro averages over three states, one of
which is rare and two of which are easy, can move substantially without
the rare state moving at all, and a reader cannot recover the per-state
behaviour from the aggregates. Cohen's $\kappa$ is included in the same
table because it is the only metric here that is directly comparable to
the inter-rater agreement figures reported for human scorers of rodent
EEG, and it is therefore the appropriate quantity for judging whether an
automated scorer is operating within the reliability of the standard it
is measured against.

% NOTE TO THE AUTHORS: once this table is filled, the following two
% sentences elsewhere in the manuscript must be rewritten to match it,
% not the other way round:
%   - Introduction: "with balanced precision, recall, and F1-scores
%     across states"
%   - Conclusion:   "The robust performance across all vigilance states,
%     particularly challenging REM sleep detection, establishes..."
% Both are claims about the contents of this table, asserted before the
% table existed. Delete them now and rewrite from the measured values.


% =====================================================================
% TABLE: confusion matrix
% =====================================================================
% REQUIRES_EXPERIMENT (protocol E1 in EXPERIMENTS.md):
%   sklearn.metrics.confusion_matrix on the same saved predictions used
%   for tab:perclass. Rows are the reference standard, columns the
%   prediction. Report raw counts, not row-normalized percentages -- the
%   counts are what let a reader recompute every other figure in this
%   paper, and a normalized matrix hides the support of the rare state.
%   ~5 minutes once E1 has run.
% =====================================================================
\begin{table}[t]
\centering
\small
\begin{tabular}{@{}lcccc@{}}
\toprule
 & \multicolumn{3}{c}{\textbf{Predicted}} & \\
\cmidrule(lr){2-4}
\textbf{Reference} & Wake & SWS & REM & \textbf{Support} \\
\midrule
Wake & \REQEXP & \REQEXP & \REQEXP & \REQEXP \\
SWS  & \REQEXP & \REQEXP & \REQEXP & \REQEXP \\
REM  & \REQEXP & \REQEXP & \REQEXP & \REQEXP \\
\bottomrule
\end{tabular}
\caption{Confusion matrix under
% REQUIRES_EXPERIMENT: name the protocol and configuration
protocol~[P\,?], in raw epoch counts. Rows are the manual reference
standard, columns the model prediction; diagonal entries are correct
classifications. Cells marked \REQEXP\ are unmeasured slots, not results.}
\label{tab:confusion}
\end{table}

% ---------------------------------------------------------------------
% PROSE FOR tab:confusion
% ---------------------------------------------------------------------
Table~\ref{tab:confusion} reports the full error structure in raw counts.
It is included because the direction of the errors, not their number,
determines whether an automated scorer is usable for a given experiment:
a scorer that confuses REM with SWS and one that confuses REM with wake
have the same macro F1 but different consequences for a study of sleep
architecture, and neither is visible in an aggregate score. Raw counts
are given rather than row-normalized rates so that every other figure in
this paper can be recomputed from this table.


% =====================================================================
% TABLE: protocol ladder
% =====================================================================
% REQUIRES_EXPERIMENT (protocols E1, E5, E6 in EXPERIMENTS.md):
%   Fill P2 from the contiguous-block split and P3 from the mean over the
%   eight leave-one-animal-out folds (mean +/- sd across folds).
%   The P1 row is the only row in this file that carries measured values,
%   and they are the figures already reported for the originally
%   submitted configuration.
%
%   IF THE P3 ROW COMES BACK BELOW THE P1 ROW -- which is the expected
%   outcome -- do not bury it, do not average the protocols together, and
%   do not report only the best fold. See EXPERIMENTS.md, experiment E6,
%   "If the result is unfavourable". The gap between these rows is the
%   most transferable finding this paper can report.
% =====================================================================
\begin{table}[t]
\centering
\footnotesize
\setlength{\tabcolsep}{3.5pt}
\begin{tabular}{@{}llcccc@{}}
\toprule
& \textbf{Held-out unit} & \textbf{Acc.} & \textbf{Macro-F1} & $\kappa$ & \textbf{REM rec.} \\
\midrule
P1 & Random epochs   & 91.5\% & 83.5\% & \REQEXP & \REQEXP \\
P2 & Time blocks     & \REQEXP & \REQEXP & \REQEXP & \REQEXP \\
P3 & Whole animals   & \REQEXP & \REQEXP & \REQEXP & \REQEXP \\
\bottomrule
\end{tabular}
\caption{The same model and hyperparameters evaluated under three
partitions of increasing strictness. The P1 accuracy and F1 entries are
the figures obtained by the originally submitted configuration; every
other cell is an unmeasured slot marked \REQEXP. P3 entries are means
over the eight folds of Table~\ref{tab:loao}. Differences between rows
are attributable to the partition alone, since nothing else changes.}
\label{tab:ladder}
\end{table}

% ---------------------------------------------------------------------
% PROSE FOR tab:ladder
% ---------------------------------------------------------------------
Table~\ref{tab:ladder} places the three protocols side by side. Its
purpose is to separate three quantities that a single accuracy figure
conflates: how well the classifier interpolates between temporally
adjacent epochs (P1 against P2), how much of its performance depends on
having seen the test animal during training (P2 against P3), and what
remains when neither source of information is available (P3 alone). Only
the last of these supports a claim about a new animal. We report the
whole ladder rather than the most favourable rung because the size of the
drop, if there is one, is the quantity a laboratory needs in order to
decide whether a scorer trained on its own cohort will work on the next
one.


% =====================================================================
% TABLE: leave-one-animal-out
% =====================================================================
% REQUIRES_EXPERIMENT (protocol E6 in EXPERIMENTS.md):
%   GroupKFold(n_splits=8) grouped on the animal identifier. Identical
%   feature pipeline, identical hyperparameters, identical seed. Fit
%   eight times. Fill 8 x 5 = 40 fold cells plus the mean and sd rows.
%   Estimated compute: 8 x (one fit) ~= 35 min at the paper's stated
%   sub-four-minute fit time; budget half a day including the plumbing.
%
%   The subject-identity column MUST be removed before this is run.
%   Under P3 the held-out animal's identifier never appears in training,
%   so the model has no fitted behaviour for it; leaving the column in
%   makes the fold result uninterpretable rather than merely pessimistic.
% =====================================================================
\begin{table}[t]
\centering
\footnotesize
\setlength{\tabcolsep}{3.5pt}
\begin{tabular}{@{}lccccc@{}}
\toprule
\textbf{Held out} & $n_{\text{test}}$ & \textbf{Acc.} & \textbf{Macro-F1} & $\kappa$ & \textbf{REM rec.} \\
\midrule
A1 & \REQEXP & \REQEXP & \REQEXP & \REQEXP & \REQEXP \\
A2 & \REQEXP & \REQEXP & \REQEXP & \REQEXP & \REQEXP \\
A3 & \REQEXP & \REQEXP & \REQEXP & \REQEXP & \REQEXP \\
A4 & \REQEXP & \REQEXP & \REQEXP & \REQEXP & \REQEXP \\
A5 & \REQEXP & \REQEXP & \REQEXP & \REQEXP & \REQEXP \\
A6 & \REQEXP & \REQEXP & \REQEXP & \REQEXP & \REQEXP \\
A7 & \REQEXP & \REQEXP & \REQEXP & \REQEXP & \REQEXP \\
A8 & \REQEXP & \REQEXP & \REQEXP & \REQEXP & \REQEXP \\
\midrule
Mean & --- & \REQEXP & \REQEXP & \REQEXP & \REQEXP \\
SD   & --- & \REQEXP & \REQEXP & \REQEXP & \REQEXP \\
Min  & --- & \REQEXP & \REQEXP & \REQEXP & \REQEXP \\
\bottomrule
\end{tabular}
\caption{Leave-one-animal-out cross-validation (protocol P3). Each row
holds out every epoch of one animal; the model is refitted from scratch
on the remaining seven with unchanged hyperparameters. Cells marked
\REQEXP\ are unmeasured slots, not results. The spread across folds, not
the mean alone, is what characterises this estimator on a cohort of
eight; the worst fold is reported for that reason.}
\label{tab:loao}
\end{table}

% ---------------------------------------------------------------------
% PROSE FOR tab:loao
% ---------------------------------------------------------------------
Table~\ref{tab:loao} reports each fold of the leave-one-animal-out
protocol separately rather than only their mean. With eight animals the
mean is an average of eight highly variable quantities, and the spread
between folds carries information the mean discards: it indicates how
much of the estimate depends on which animals happen to be in the
training set, and it is the closest thing this cohort can offer to an
estimate of what would happen on a ninth animal. The minimum row is
included because a laboratory adopting an automated scorer needs the
worst case it is likely to encounter, not the average case. Per-fold
variation in the class prior, given in Table~\ref{tab:classdist}, should
be read alongside these rows.

% NOTE TO THE AUTHORS: eight paired folds is also what makes a
% significance test possible at all. A Wilcoxon signed-rank test over
% these eight folds is the only place in this paper where the word
% "significantly" can be used honestly; it must not be used anywhere
% else until that test has been run. See EXPERIMENTS.md, experiment E10.


% =====================================================================
% TABLE: feature-family ablation
% =====================================================================
% REQUIRES_EXPERIMENT (protocol E9 in EXPERIMENTS.md):
%   Two blocks. Cumulative addition establishes what each family adds on
%   top of the previous ones; leave-one-family-out establishes what each
%   family contributes in the presence of all the others. Both are needed
%   because the families are strongly collinear -- Hjorth parameters are
%   functions of spectral moments, and an amplitude-range statistic
%   correlates with delta power -- so a family can look essential in one
%   block and redundant in the other. That divergence is itself a result
%   worth reporting.
%   Run every variant under P3 on all eight folds with >= 5 seeds and
%   report mean +/- sd. Estimated compute: 8 variants x 8 folds x 5 seeds
%   = 320 fits; at the paper's stated sub-four-minute fit time that is
%   roughly 21 h of compute. Budget one day wall clock.
%
%   THE CFC ROWS PRESUPPOSE THAT A CROSS-FREQUENCY COUPLING FEATURE
%   EXISTS IN THE CODE. If it does not, these two rows must be deleted
%   rather than filled, and the discrepancy with the title and the frozen
%   abstract handled separately. Do not construct a coupling feature
%   after the fact and present it as the one the abstract described.
% =====================================================================
\begin{table}[t]
\centering
\footnotesize
\setlength{\tabcolsep}{4pt}
\begin{tabular}{@{}lcccc@{}}
\toprule
\textbf{Feature set} & $d$ & \textbf{Acc.} & \textbf{Macro-F1} & \textbf{REM rec.} \\
\midrule
\multicolumn{5}{@{}l}{\emph{Cumulative addition}}\\
Spectral only        & \REQEXP & \REQEXP & \REQEXP & \REQEXP \\
\quad + Hjorth       & \REQEXP & \REQEXP & \REQEXP & \REQEXP \\
\quad + MMD          & \REQEXP & \REQEXP & \REQEXP & \REQEXP \\
\quad + CFC (full)   & \REQEXP & \REQEXP & \REQEXP & \REQEXP \\
\addlinespace
\multicolumn{5}{@{}l}{\emph{Leave-one-family-out from the full set}}\\
Full $-$ Spectral    & \REQEXP & \REQEXP & \REQEXP & \REQEXP \\
Full $-$ Hjorth      & \REQEXP & \REQEXP & \REQEXP & \REQEXP \\
Full $-$ MMD         & \REQEXP & \REQEXP & \REQEXP & \REQEXP \\
Full $-$ CFC         & \REQEXP & \REQEXP & \REQEXP & \REQEXP \\
\bottomrule
\end{tabular}
\caption{Feature-family ablation under protocol P3, mean over the eight
folds and over
% REQUIRES_EXPERIMENT: number of seeds
seeds. $d$ is the dimensionality of the design matrix for that variant.
Cells marked \REQEXP\ are unmeasured slots, not results. The two blocks
answer different questions and need not agree: the upper block measures
what a family adds when little else is present, the lower what it adds
when everything else is present.}
\label{tab:ablation}
\end{table}

% ---------------------------------------------------------------------
% PROSE FOR tab:ablation
% ---------------------------------------------------------------------
Table~\ref{tab:ablation} removes each feature family and refits. This is
reported in addition to the attribution analysis because the two measure
different things. A gain-based importance score is a property of one
fitted model and is known to distribute credit arbitrarily among
correlated predictors, and the families used here are correlated by
construction: the Hjorth descriptors are functions of the second and
fourth spectral moments and therefore covary with the band powers, and an
amplitude-range statistic covaries with low-frequency power. An ablation
answers the question the attribution analysis cannot---whether the model
can do without a family---and it is the only evidence in this paper that
bears on whether any individual descriptor is doing work.

% NOTE TO THE AUTHORS: the manuscript's Conclusion currently asserts that
% "The MMD feature successfully captured unique temporal dynamics of
% neural oscillations, enhancing discrimination between..." That is a
% claim about the "Full - MMD" row of this table, made before the row
% existed. Delete it now. When the row is measured, write whatever it
% says. A negative result on MMD is entirely publishable -- the frozen
% abstract lists MMD as a feature that was used, which remains true
% whatever the ablation shows -- and it is far less damaging than a
% reviewer inferring that the feature was never tested.


% =====================================================================
% TABLE: baseline comparison
% =====================================================================
% REQUIRES_EXPERIMENT (protocols E7, E11 in EXPERIMENTS.md):
%   Every row must be measured on THESE eight recordings against THESE
%   labels under the stated protocol. No row may be transcribed from
%   another paper -- external figures belong in tab:context, which is
%   explicitly labelled as not comparable.
%   The majority-class row follows from tab:classdist alone.
% =====================================================================
\begin{table}[t]
\centering
\footnotesize
\setlength{\tabcolsep}{3pt}
\begin{tabular}{@{}lcccc@{}}
\toprule
\textbf{Method} & \textbf{Acc.} & \textbf{Mac.-F1} & $\kappa$ & \textbf{REM rec.} \\
\midrule
Majority class            & \REQEXP & \REQEXP & \REQEXP & \REQEXP \\
Logistic regression       & \REQEXP & \REQEXP & \REQEXP & \REQEXP \\
Feed-forward network      & \REQEXP & \REQEXP & \REQEXP & \REQEXP \\
SAGER$^{\dagger}$         & \REQEXP & \REQEXP & \REQEXP & \REQEXP \\
\addlinespace
XGBoost (this work)       & \REQEXP & \REQEXP & \REQEXP & \REQEXP \\
\quad + temporal smooth.  & \REQEXP & \REQEXP & \REQEXP & \REQEXP \\
\bottomrule
\end{tabular}
\caption{Baseline comparison, all methods evaluated on the same eight
recordings against the same reference labels under protocol
% REQUIRES_EXPERIMENT: name the protocol; produce one instance per
% protocol or add a protocol column
[P\,?]. Cells marked \REQEXP\ are unmeasured slots, not results.
$^{\dagger}$SAGER \citep{saevskiy2025_sensors} is an external, published,
open-source single-channel rodent scorer and is the only row not
implemented by us. No figure in this table is taken from another
publication; figures reported elsewhere on other cohorts are listed
separately in Table~\ref{tab:context} and are not comparable to these.}
\label{tab:baselines}
\end{table}

% ---------------------------------------------------------------------
% PROSE FOR tab:baselines
% ---------------------------------------------------------------------
Table~\ref{tab:baselines} compares the proposed classifier against a
trivial floor, two in-house alternatives, and one external published
system. The majority-class row is included because accuracy on a
three-class problem with the prior of Table~\ref{tab:classdist} is
uninterpretable without it. The external comparison matters more than the
in-house ones: a comparison against models written by the same authors
measures implementation choices as much as method, whereas SAGER
\citep{saevskiy2025_sensors} is an independently developed,
independently published, openly available scorer for exactly this
setting---single-channel rodent EEG without EMG---and running it on these
recordings is the only way to obtain a comparison that is not internal to
this work. Systems requiring a paired EMG channel, including AccuSleep
\citep{barger2019_accusleep} and IntelliSleepScorer
\citep{wang2023_intellisleepscorer}, cannot be evaluated on this cohort
at all, and we make no claim relative to them.


% =====================================================================
% TABLE: calibration  -- OPTIONAL, cut third if space binds
% =====================================================================
% REQUIRES_EXPERIMENT (protocol E8 in EXPERIMENTS.md):
%   Reliability analysis, ECE, MCE, and Brier score under each protocol,
%   and per class. Requires a probabilistic objective; see the note in
%   the hyperparameter description above.
%   Calibration is the property most likely to degrade under subject
%   shift, which is exactly why it must be measured under P3 and not only
%   under P1. If it holds under P3, that is the strongest result this
%   paper can report. If it does not, report the degradation -- a
%   measured statement that confidence scores are unreliable under
%   subject shift is a useful finding and is the honest version of the
%   closed-loop argument.
% =====================================================================
\begin{table}[t]
\centering
\footnotesize
\setlength{\tabcolsep}{4pt}
\begin{tabular}{@{}lcccc@{}}
\toprule
\textbf{Protocol} & \textbf{ECE} & \textbf{MCE} & \textbf{Brier} & \textbf{Worst class} \\
\midrule
P1 & \REQEXP & \REQEXP & \REQEXP & \REQEXP \\
P2 & \REQEXP & \REQEXP & \REQEXP & \REQEXP \\
P3 & \REQEXP & \REQEXP & \REQEXP & \REQEXP \\
\bottomrule
\end{tabular}
\caption{Calibration of the predicted class probabilities under each
protocol, computed as defined in
Eqs.~(\ref{eq:brier})--(\ref{eq:ece}). ``Worst class'' names the state
with the largest one-vs-rest deviation from the diagonal. Cells marked
\REQEXP\ are unmeasured slots, not results.}
\label{tab:calib}
\end{table}

% ---------------------------------------------------------------------
% PROSE FOR tab:calib
% ---------------------------------------------------------------------
Table~\ref{tab:calib} reports calibration under each protocol. It is
separated from the accuracy tables because the two properties can move
independently, and because calibration is the quantity that closed-loop
protocols actually depend on: an experiment that triggers a stimulus when
the predicted probability of REM exceeds a threshold is governed by
whether that probability means what it says, not by the model's accuracy
at its own arg-max. Calibration is also the property most exposed to
subject shift, which is why it is reported under P3 and not only under
P1.

% NOTE TO THE AUTHORS: the manuscript currently states that "the maximum
% deviation from the diagonal never exceeds 3 percentage points." The
% reliability diagram printed directly beneath that sentence does not
% support it. Delete the sentence now, before this table is filled, and
% write the replacement from the arrays that produced the plot -- not
% from the plot. Also relabel the figure legend from Class 0/1/2 to
% Wake/SWS/REM and state the LabelEncoder mapping, so that a reader can
% tell whether the poorly calibrated class is the minority one.


% =====================================================================
% TABLE: compute -- OPTIONAL, cut second if space binds
% =====================================================================
% REQUIRES_EXPERIMENT (protocol E13 in EXPERIMENTS.md):
%   Feature-extraction and inference throughput on a workstation CPU and
%   on an embedded-class device. Batch and single-epoch latency both, and
%   include feature extraction -- for this pipeline it is likely to
%   dominate inference, and reporting tree-traversal time alone would
%   overstate the deployability advantage.
% =====================================================================
\begin{table}[t]
\centering
\footnotesize
\setlength{\tabcolsep}{4pt}
\begin{tabular}{@{}lccc@{}}
\toprule
\textbf{Device} & \textbf{Feat.\ extr.} & \textbf{Inference} & \textbf{Total / epoch} \\
\midrule
Workstation CPU & \REQEXP & \REQEXP & \REQEXP \\
Embedded device & \REQEXP & \REQEXP & \REQEXP \\
\bottomrule
\end{tabular}
\caption{Per-epoch processing time, feature extraction and model
inference reported separately. Cells marked \REQEXP\ are unmeasured
slots, not results. Real-time operation on $10$\,s epochs requires the
total to remain below $10$\,s per epoch by a margin.}
\label{tab:compute}
\end{table}

% ---------------------------------------------------------------------
% PROSE FOR tab:compute
% ---------------------------------------------------------------------
Table~\ref{tab:compute} reports the cost of scoring one epoch end to end.
Feature extraction and inference are separated because they scale
differently and because, for a pipeline built on a few hundred shallow
trees over a low-dimensional descriptor, the spectral transform is likely
to dominate; a figure that counted only tree traversal would overstate
what the pipeline can do on constrained hardware. The relevant threshold
is set by the epoch length: untethered closed-loop use requires the total
to stay below the duration of the epoch being scored.


% =====================================================================
% TABLE: external context -- OPTIONAL, CUT FIRST if space binds
% =====================================================================
% !!! READ THIS BEFORE TOUCHING THE TABLE BELOW !!!
%
% This is the ONLY table in this file containing measured numbers, and
% NONE of them was measured on this cohort. Each is a figure reported by
% the cited authors on their own data, under their own protocol, with
% their own labels and in some cases their own species and signal set.
% They are NOT comparable to anything in tab:baselines and the caption
% says so.
%
% Provenance of each figure, for re-verification before submission:
%   IntelliSleepScorer 95.2% acc / kappa 0.91  -- Wang et al., Sci Rep
%     13, 2023, doi 10.1038/s41598-023-31288-2. EEG+EMG, mice.
%   SlumberNet 97% acc / 96% F1 -- Jha, Valekunja & Reddy, Sci Rep 14:4797,
%     2024, doi 10.1038/s41598-024-54727-0. EEG+EMG, mice.
%   Smith et al. mean F1 87.6% -- NPP Digital Psychiatry & Neuroscience,
%     2025, doi 10.1038/s44277-025-00035-y. Single-channel EEG, rats.
%   SAGER 92% (rats) / 93% (mice) -- Saevskiy et al., Sensors 25(3):921,
%     2025, doi 10.3390/s25030921. NOTE: these are SLEEP-WAKE detection
%     accuracies, not three-state accuracies. Do not present them as
%     three-state accuracy.
%
% ACTION REQUIRED: open each publisher record and confirm each figure
% before this table appears in a submitted PDF. Six of the seven domain
% citations originally in this manuscript were wrong; one more bad
% number here is unrecoverable.
%
% If in doubt, DELETE THIS TABLE. It is context, not evidence, and the
% same content can be carried in one Related Work sentence per system.
% =====================================================================
\begin{table}[t]
\centering
\footnotesize
\setlength{\tabcolsep}{4pt}
\begin{tabular}{@{}llc@{}}
\toprule
\textbf{System} & \textbf{Signals} & \textbf{Reported} \\
\midrule
IntelliSleepScorer \citep{wang2023_intellisleepscorer} & EEG+EMG & 95.2\% acc.$^{\ddagger}$ \\
SlumberNet \citep{jha2024_slumbernet}                  & EEG+EMG & 97\% acc.$^{\ddagger}$ \\
\citet{smith2025_ratlstm}                              & EEG     & 87.6\% F1$^{\ddagger}$ \\
SAGER \citep{saevskiy2025_sensors}                     & EEG     & 92--93\% acc.$^{\ddagger\S}$ \\
\bottomrule
\end{tabular}
\caption{Figures reported by other authors for other rodent scorers.
$^{\ddagger}$Measured by the cited authors on their own cohorts, under
their own protocols and against their own reference labels; these are
\emph{not} comparable to Table~\ref{tab:baselines} and no claim in this
paper is made relative to them. $^{\S}$Sleep--wake detection accuracy,
not three-state accuracy. The table is included to indicate where this
application area currently sits, not to position this work within it.}
\label{tab:context}
\end{table}

% ---------------------------------------------------------------------
% PROSE FOR tab:context
% ---------------------------------------------------------------------
Table~\ref{tab:context} lists figures published for other rodent scorers.
It is included as context for the reader and not as a comparison: the
entries differ from this work and from each other in species, in signal
set, in cohort size, in evaluation protocol, and in the reference
standard used, and no pair of numbers in the table was produced under
conditions that would make their difference interpretable. The only
comparison in this paper that controls those factors is
Table~\ref{tab:baselines}, in which every method is run on the same
recordings against the same labels.

% =====================================================================
% END OF FILE
%
% BIBLIOGRAPHY DEPENDENCY. This file cites the following keys. Those
% marked (new) do not exist in the current references.bib and must be
% added; the verified entries are in the revision plan (item A1) and in
% new_refs.bib, except barger2019_accusleep which is given below.
%
%   Chen_2016                    -- present
%   saevskiy2025_sensors         -- present (verify the title: the bib
%                                   abbreviates "Electroencephalogram"
%                                   to EEG; the published title spells
%                                   it out)
%   e18090272                    -- present
%   wang2023_intellisleepscorer  -- (new) revision plan A1.7
%   jha2024_slumbernet           -- (new) revision plan A1.2
%   smith2025_ratlstm            -- (new) revision plan A1.4
%   barger2019_accusleep         -- (new) entry below
%
% Verified this session against the AccuSleep repository README and the
% PLOS ONE record:
%
% @article{barger2019_accusleep,
%   author  = {Barger, Zeke and Frye, Charles G. and Liu, Danqian and
%              Dan, Yang and Bouchard, Kristofer E.},
%   title   = {Robust, Automated Sleep Scoring by a Compact Neural
%              Network with Distributional Shift Correction},
%   journal = {PLOS ONE},
%   volume  = {14},
%   number  = {12},
%   pages   = {e0224642},
%   year    = {2019},
%   doi     = {10.1371/journal.pone.0224642}
% }
%
% NOT CITED HERE, deliberately: no citation is made for the claim that
% no published rodent scorer reports probability calibration. If that
% claim is used anywhere in the paper it must first be checked against
% each cited system individually. It is a strong negative claim and it
% must survive one search by one reviewer.
% =====================================================================
